% GENERATED FILE. Edit canonical lessons, then run npm run generate:books.
% canonical-source-hash: fnv1a64:4b7f256bd4d62193
% canonical-lessons: JA-C415-yuudachi, JA-C415-kosame, JA-C415-bisshori, JA-C415-nurasu, JA-C415-amayadori

\chapter{Caught in a Shower}
\label{ch:ja-caught-in-a-shower}
\hlchaptermodality{415}

\begin{chapteropening}
\textbf{By the end of this chapter:} \emph{I can talk about being caught in the rain.}

Complete the last lesson of chapter 415: I can talk about being caught in the rain.
\end{chapteropening}

\section[\texorpdfstring{\ja{ゆうだち} (yūdachi)}{ゆうだち (yūdachi)}]{\ja{ゆうだち} (yūdachi) — a sudden evening shower}

\label{lesson:JA-C415-yuudachi}



\begin{quote}
Before the new one: say the Japanese for shade, then the Japanese for dazzling.
\end{quote}

\subsection*{You'll want to know: \ja{ゆうだち}}
\textbf{\ja{ゆうだち}} — \emph{yūdachi} — \textquotedblleft{}a sudden evening shower\textquotedblright{}.

\begin{quote}
\textbf{\ja{なつの} \ja{ごごに} \ja{ゆうだちが} \ja{ふりました。}} — \emph{natsu no gogo ni yūdachi ga furimashita} — \textquotedblleft{}A sudden shower fell on the summer afternoon.\textquotedblright{}
\end{quote}

\begin{cousinweb}
\textbf{\ja{ゆうだち}} (\emph{yūdachi}) is a native Japanese word (wago), inherited rather than borrowed from Chinese, built from \emph{yū}, \textquotedblleft{}evening\textquotedblright{} and \emph{dachi}, \textquotedblleft{}rising\textquotedblright{} (from \emph{tatsu}); its kanji were attached later as kun'yomi (JMdict).
\end{cousinweb}

\begin{grammarlens}[title={What you've built}]
One more word for getting around.
\end{grammarlens}

\subsection*{Your turn}
\begin{itemize}
\raggedright
  \item \emph{Say it:} \emph{yūdachi}
  \item \emph{Say it:} \emph{yūdachi}, once more
  \item \emph{Say it:} \emph{natsu no gogo ni yūdachi ga furimashita} — A sudden shower fell on the summer afternoon.
  \item \emph{Recall:} say the Japanese for shade, then the Japanese for dazzling, then say \emph{yūdachi} again
\end{itemize}

\begin{tcolorbox}[breakable,colback=teal!4,colframe=teal!35!black,title={Before you move on}]
What does \ja{ゆうだち} mean? (\textquotedblleft{}A sudden evening shower\textquotedblright{}.) Say it once more.
\end{tcolorbox}
\section[\texorpdfstring{\ja{こさめ} (kosame)}{こさめ (kosame)}]{\ja{こさめ} (kosame) — light rain; drizzle}

\label{lesson:JA-C415-kosame}



\begin{quote}
Before the new one: say the Japanese for dazzling, then the Japanese for a sudden evening shower.
\end{quote}

\subsection*{You'll want to know: \ja{こさめ}}
\textbf{\ja{こさめ}} — \emph{kosame} — \textquotedblleft{}light rain; drizzle\textquotedblright{}.

\begin{quote}
\textbf{\ja{こさめの} \ja{なかを} \ja{あるきました。}} — \emph{kosame no naka o arukimashita} — \textquotedblleft{}I walked in the light rain.\textquotedblright{}
\end{quote}

\begin{cousinweb}
\textbf{\ja{こさめ}} (\emph{kosame}) is a native Japanese word (wago), inherited rather than borrowed from Chinese, built from \emph{ko}, \textquotedblleft{}small\textquotedblright{} and \emph{same}, \textquotedblleft{}rain\textquotedblright{} (from \emph{ame}); its kanji were attached later as kun'yomi (JMdict).
\end{cousinweb}

\begin{grammarlens}[title={What you've built}]
One more word for getting around.
\end{grammarlens}

\subsection*{Your turn}
\begin{itemize}
\raggedright
  \item \emph{Say it:} \emph{kosame}
  \item \emph{Say it:} \emph{kosame}, once more
  \item \emph{Say it:} \emph{kosame no naka o arukimashita} — I walked in the light rain.
  \item \emph{Recall:} say the Japanese for dazzling, then the Japanese for a sudden evening shower, then say \emph{kosame} again
\end{itemize}

\begin{tcolorbox}[breakable,colback=teal!4,colframe=teal!35!black,title={Before you move on}]
What does \ja{こさめ} mean? (\textquotedblleft{}Light rain; drizzle\textquotedblright{}.) Say it once more.
\end{tcolorbox}
\section[\texorpdfstring{\ja{びっしょり} (bisshori)}{びっしょり (bisshori)}]{\ja{びっしょり} (bisshori) — soaking wet}

\label{lesson:JA-C415-bisshori}



\begin{quote}
Before the new one: say the Japanese for a sudden evening shower, then the Japanese for light rain.
\end{quote}

\subsection*{You'll want to know: \ja{びっしょり}}
\textbf{\ja{びっしょり}} — \emph{bisshori} — \textquotedblleft{}soaking wet\textquotedblright{}.

\begin{quote}
\textbf{\ja{あめで} \ja{ふくが} \ja{びっしょり} \ja{ぬれました。}} — \emph{ame de fuku ga bisshori nuremashita} — \textquotedblleft{}My clothes got soaking wet in the rain.\textquotedblright{}
\end{quote}

\begin{cousinweb}
\textbf{\ja{びっしょり}} (\emph{bisshori}) is a native mimetic word that pictures being drenched right through: an imitative Japanese form, inherited rather than borrowed, whose older origin is not settled (JMdict).
\end{cousinweb}

\begin{grammarlens}[title={What you've built}]
One more word for getting around.
\end{grammarlens}

\subsection*{Your turn}
\begin{itemize}
\raggedright
  \item \emph{Say it:} \emph{bisshori}
  \item \emph{Say it:} \emph{bisshori}, once more
  \item \emph{Say it:} \emph{ame de fuku ga bisshori nuremashita} — My clothes got soaking wet in the rain.
  \item \emph{Recall:} say the Japanese for a sudden evening shower, then the Japanese for light rain, then say \emph{bisshori} again
\end{itemize}

\begin{tcolorbox}[breakable,colback=teal!4,colframe=teal!35!black,title={Before you move on}]
What does \ja{びっしょり} mean? (\textquotedblleft{}Soaking wet\textquotedblright{}.) Say it once more.
\end{tcolorbox}
\section[\texorpdfstring{\ja{ぬらす} (nurasu)}{ぬらす (nurasu)}]{\ja{ぬらす} (nurasu) — to get something wet; to wet}

\label{lesson:JA-C415-nurasu}



\begin{quote}
Before the new one: say the Japanese for light rain, then the Japanese for soaking wet.
\end{quote}

\subsection*{You'll want to know: \ja{ぬらす}}
\textbf{\ja{ぬらす}} — \emph{nurasu} — \textquotedblleft{}to get something wet; to wet\textquotedblright{}.

\begin{quote}
\textbf{\ja{あめで} \ja{くつを} \ja{ぬらしました。}} — \emph{ame de kutsu o nurashimashita} — \textquotedblleft{}I got my shoes wet in the rain.\textquotedblright{}
\end{quote}

Polite \textbf{\ja{ぬらします}}, \emph{nurashimasu}; a godan verb. Past (polite): \textbf{\ja{こどもが} \ja{ふくを} \ja{ぬらしました。}}, \emph{kodomo ga fuku o nurashimashita}.

\begin{cousinweb}
\textbf{\ja{ぬらす}} (\emph{nurasu}) is a native Japanese word (wago), inherited rather than borrowed from Chinese; its kanji was attached later as a kun'yomi (JMdict). It is the partner of \emph{nureru}, \textquotedblleft{}to get wet\textquotedblright{} (JMdict).
\end{cousinweb}

\begin{grammarlens}[title={What you've built}]
One more word for getting around.
\end{grammarlens}

\subsection*{Your turn}
\begin{itemize}
\raggedright
  \item \emph{Say it:} \emph{nurasu}
  \item \emph{Say it:} \emph{nurasu}, once more
  \item \emph{Say it:} \emph{ame de kutsu o nurashimashita} — I got my shoes wet in the rain.
  \item \emph{Recall:} say the Japanese for light rain, then the Japanese for soaking wet, then say \emph{nurasu} again
\end{itemize}

\begin{tcolorbox}[breakable,colback=teal!4,colframe=teal!35!black,title={Before you move on}]
What does \ja{ぬらす} mean? (\textquotedblleft{}To get something wet; to wet\textquotedblright{}.) Say it once more.
\end{tcolorbox}
\section[\texorpdfstring{\ja{あまやどり} (amayadori)}{あまやどり (amayadori)}]{\ja{あまやどり} (amayadori) — sheltering from the rain}

\label{lesson:JA-C415-amayadori}



\begin{quote}
Before the new one: say the Japanese for soaking wet, then the Japanese for to get something wet.
\end{quote}

\subsection*{You'll want to know: \ja{あまやどり}}
\textbf{\ja{あまやどり}} — \emph{amayadori} — \textquotedblleft{}sheltering from the rain\textquotedblright{}.

\begin{quote}
\textbf{\ja{えきで} \ja{あまやどりを} \ja{しました。}} — \emph{eki de amayadori o shimashita} — \textquotedblleft{}I sheltered from the rain at the station.\textquotedblright{}
\end{quote}

\begin{cousinweb}
\textbf{\ja{あまやどり}} (\emph{amayadori}) is a native Japanese word (wago), inherited rather than borrowed from Chinese, built from \emph{ama}, \textquotedblleft{}rain\textquotedblright{} (from \emph{ame}) and \emph{yadori}, \textquotedblleft{}lodging\textquotedblright{}; its kanji were attached later as kun'yomi (JMdict).
\end{cousinweb}

\begin{grammarlens}[title={What you've built}]
One more word for getting around.
\end{grammarlens}

\subsection*{Your turn}
\begin{itemize}
\raggedright
  \item \emph{Say it:} \emph{amayadori}
  \item \emph{Say it:} \emph{amayadori}, once more
  \item \emph{Say it:} \emph{eki de amayadori o shimashita} — I sheltered from the rain at the station.
  \item \emph{Recall:} say the Japanese for soaking wet, then the Japanese for to get something wet, then say \emph{amayadori} again
\end{itemize}

\begin{tcolorbox}[breakable,colback=teal!4,colframe=teal!35!black,title={Before you move on}]
What does \ja{あまやどり} mean? (\textquotedblleft{}Sheltering from the rain\textquotedblright{}.) Say it once more.
\end{tcolorbox}
